\documentclass[11pt]{article}

\usepackage{amsmath,amssymb,amsthm,mathtools}
\usepackage[hidelinks]{hyperref}
\usepackage{microtype}

\title{ExpByOpt: the $\Lambda$-game, minimax duality, and the information-ratio bound}
\date{}

% --- basic notation ---
\newcommand{\R}{\mathbb{R}}
\newcommand{\Rd}{\mathbb{R}^d}
\newcommand{\E}{\mathbb{E}}
\newcommand{\ip}[2]{\langle #1,#2\rangle}
\newcommand{\cA}{\mathcal{A}}
\newcommand{\cL}{\mathcal{L}}
\newcommand{\cP}{\mathcal{P}}
\newcommand{\pr}{\bar A_\star}   % prior mean
\newcommand{\po}{\Pi}           % posterior mean (random)

% --- theorem environments ---
\theoremstyle{definition}
\newtheorem{definition}{Definition}
\newtheorem{assumption}{Assumption}
\theoremstyle{plain}
\newtheorem{lemma}{Lemma}
\newtheorem{theorem}{Theorem}
\theoremstyle{remark}
\newtheorem{remark}{Remark}
\newtheorem{conjecture}{Conjecture}

\begin{document}
\maketitle

\section{Setup (only what is needed for $\Lambda$)}

\begin{assumption}[Bandit linear feedback]
    Fix a convex compact action set $\cA\subset \Rd$ and a bounded loss set $\cL\subset \Rd$.
    In one round, an adversary chooses $(a_\star,\ell)\in \cA\times \cL$ and the learner:
    (i) chooses a distribution $P\in \cP(\cA)$, samples $A\sim P$,
    (ii) observes the scalar $Y=\ip{A}{\ell}$,
    (iii) outputs an estimate $\hat\ell = E(A,Y)\in\Rd$, where $E:\cA\times\R\to\Rd$ is an estimator.
\end{assumption}

\begin{assumption}[Mirror map]
    Fix a Legendre potential $F:\Rd\to\R\cup\{+\infty\}$ with $\cA\subset \mathrm{int}(\mathrm{dom}(F))$.
    Let $F^\star$ be its convex conjugate. Define the Bregman divergences
    \[
        D_F(x,y)=F(x)-F(y)-\ip{\nabla F(y)}{x-y},
        \qquad
        D_{F^\star}(u,v)=F^\star(u)-F^\star(v)-\ip{\nabla F^\star(v)}{u-v}.
    \]
    We use the standard identities:
    \[
        D_F(x,y)=D_{F^\star}(\nabla F(y),\nabla F(x)),\qquad
        \ip{\nabla F(z)-\nabla F(y)}{x-y}=D_F(x,y)+D_F(y,z)-D_F(x,z).
    \]
\end{assumption}

\section{The per-round functional $\Lambda$ and why ExpByOpt cares about it}

Fix a reference point $a\in\cA$ (think: the ``base'' action suggested by the full-information mirror descent update) and a stepsize $\eta>0$.
Given a sampling rule $P$ and estimator $E$, define $A\sim P$, $Y=\ip{A}{\ell}$, and $\hat\ell=E(A,Y)$.

\begin{definition}[$\Lambda$ functional]
    For $(a_\star,\ell)\in\cA\times\cL$ define
    \begin{equation}
        \label{eq:Lambda-point}
        \Lambda_{a,\eta}(P,E \,\Vert\, a_\star,\ell)
        \;=\;
        \frac{1}{\eta}\,
        \E\!\left[
        \ip{A-a_\star}{\ell}
        -\ip{a-a_\star}{\hat\ell}
        +\frac{1}{\eta}\,D_{F^\star}\!\big(\nabla F(a)-\eta \hat\ell,\nabla F(a)\big)
        \right],
    \end{equation}
    where the expectation is over $A\sim P$ (hence over $(Y,\hat\ell)$).
\end{definition}

The quantity in \eqref{eq:Lambda-point} is the single-round term that appears in the standard regret analysis of \emph{biased stochastic mirror descent} (see, e.g., the bandit mirror descent / ExpByOpt development in Lattimore--Szepesv\'ari).

\begin{definition}[Worst-case per-round value at fixed $\eta$]
    Define
    \[
        V_\eta(a)
        \;=\;
        \inf_{P\in\cP(\cA)}\ \inf_{E:\cA\times\R\to\Rd}\ \sup_{(a_\star,\ell)\in\cA\times\cL}
        \Lambda_{a,\eta}(P,E \,\Vert\, a_\star,\ell),
        \qquad
        \Lambda^\star_\eta
        \;=\;
        \sup_{a\in\cA} V_\eta(a).
    \]
\end{definition}

\begin{remark}[How this drives regret]
    ExpByOpt can be viewed as: at each round, given the current base point $a_t$, choose $(P_t,E_t)$ that approximately attains $V_\eta(a_t)$.
    Then the standard analysis gives an $n$-round regret bound of the form
    \[
        \E[\mathrm{Reg}_n(a_\star)]
        \;\le\;
        \frac{D_F(a_\star,a_1)}{\eta}
        \;+\;
        n\,\eta\,(\Lambda^\star_\eta+\text{(optimization error)}).
    \]
    All the structure below is about upper bounding $\Lambda^\star_\eta$ via minimax duality and connecting it to an information ratio.
\end{remark}

\section{Randomized adversary and minimax duality}

To make the adversary strategy set convex, let it randomize $(A_\star,\ell)$.
Let $\nu\in\cP(\cA\times\cL)$ and sample $(A_\star,\ell)\sim \nu$ independently of $A\sim P$.

\begin{definition}[$\Lambda$ against a distribution $\nu$]
    \begin{equation}
        \label{eq:Lambda-nu}
        \Lambda_{a,\eta}(P,E \,\Vert\, \nu)
        \;=\;
        \frac{1}{\eta}\,
        \E_{(A_\star,\ell)\sim\nu,\ A\sim P}\!\left[
        \ip{A-A_\star}{\ell}
        -\ip{a-A_\star}{\hat\ell}
        +\frac{1}{\eta}\,D_{F^\star}\!\big(\nabla F(a)-\eta \hat\ell,\nabla F(a)\big)
        \right].
    \end{equation}
\end{definition}

Clearly,
\[
    \sup_{(a_\star,\ell)\in\cA\times\cL}\Lambda_{a,\eta}(P,E\Vert a_\star,\ell)
    \;\le\;
    \sup_{\nu\in\cP(\cA\times\cL)}\Lambda_{a,\eta}(P,E\Vert \nu).
\]

\begin{lemma}[Minimax swap (Sion-type)]
    \label{lem:minimax}
    For fixed $a\in\cA$ and $\eta>0$, under standard compactness/continuity conditions,
    \[
        \inf_{P}\ \inf_{E}\ \sup_{\nu}\ \Lambda_{a,\eta}(P,E\Vert \nu)
        \;=\;
        \sup_{\nu}\ \inf_{P}\ \inf_{E}\ \Lambda_{a,\eta}(P,E\Vert \nu).
    \]
\end{lemma}

This is the ``duality'' step used to analyze the game via a (Bayesian) information quantity.

\section{Optimal unrestricted estimator and the information-ratio bound}

Fix $(a,\eta,P,\nu)$ and minimize \eqref{eq:Lambda-nu} over all estimators $E:\cA\times\R\to\Rd$.

Define the \emph{prior mean} and \emph{posterior mean}:
\[
    \pr \;=\; \E_\nu[A_\star],
    \qquad
    \po \;=\; \E_\nu[A_\star \mid A,Y],
\]
where randomness is over $(A_\star,\ell)\sim\nu$ and $A\sim P$ (so $\po$ is a random vector in $\cA$).

Also define the (Bayesian) instantaneous regret gap
\begin{equation}
    \label{eq:Delta}
    \Delta(\nu,P)
    \;=\;
    \E_{(A_\star,\ell)\sim\nu,\ A\sim P}\!\big[\ip{A-A_\star}{\ell}\big].
\end{equation}

\begin{lemma}[Optimal unrestricted estimator]
    \label{lem:opt-unrestricted}
    For fixed $(a,\eta,P,\nu)$, an optimal estimator is given pointwise by
    \[
        G^\nu(b,y)
        \;=\;
        \frac{1}{\eta}\Big(\nabla F(a)-\nabla F\big(\E_\nu[A_\star\mid A=b,\ Y=y]\big)\Big),
    \]
    and the resulting value satisfies
    \begin{equation}
        \label{eq:Lambda-unrestricted}
        \Lambda_{a,\eta}(P,G^\nu \Vert \nu)
        \;=\;
        \frac{1}{\eta}\,\Delta(\nu,P)
        \;-\;
        \frac{1}{\eta^2}\,\E\!\big[D_F(\po,a)\big].
    \end{equation}
\end{lemma}

\begin{lemma}[Expected Bregman identity]
    \label{lem:EBI}
    If $F$ is strictly convex and differentiable, then for any random vector $X$,
    \[
        \E[D_F(X,a)]
        \;=\;
        \E[D_F(X,\E[X])]
        \;+\;
        D_F(\E[X],a).
    \]
    In particular, since $\E[\po]=\pr$,
    \begin{equation}
        \label{eq:EBI-po}
        \E[D_F(\po,a)]
        \;=\;
        \E[D_F(\po,\pr)]
        \;+\;
        D_F(\pr,a).
    \end{equation}
\end{lemma}

Combining \eqref{eq:Lambda-unrestricted} and \eqref{eq:EBI-po} and dropping $D_F(\pr,a)\ge 0$ yields
\begin{equation}
    \label{eq:Lambda-unrestricted-drop}
    \Lambda_{a,\eta}(P,G^\nu \Vert \nu)
    \;\le\;
    \frac{1}{\eta}\,\Delta(\nu,P)
    \;-\;
    \frac{1}{\eta^2}\,\underbrace{\E[D_F(\po,\pr)]}_{\text{``information gain''}}.
\end{equation}

\begin{theorem}[Information-ratio upper bound (general Bregman form)]
    \label{thm:info-ratio}
    Let $I_F(\nu,P)=\E[D_F(\po,\pr)]$. If $I_F(\nu,P)>0$, then for any $\eta>0$,
    \[
        \Lambda_{a,\eta}(P,G^\nu \Vert \nu)
        \;\le\;
        \frac{\Delta(\nu,P)^2}{4\,I_F(\nu,P)}.
    \]
\end{theorem}

\begin{remark}[Connection to the classical information ratio]
    If $\cA=\Delta(k)$, $F$ is (negative) entropy so $D_F=\mathrm{KL}$, and $A_\star$ is a one-hot vector encoding the identity of the optimal arm,
    then $\po$ is exactly the posterior distribution of $A_\star$ and $\pr$ is its prior distribution. In that case,
    $I_F(\nu,P)=\E[\mathrm{KL}(\po\Vert \pr)]$ equals the mutual information $I(A_\star;(A,Y))$,
    and $\Delta(\nu,P)^2/I_F(\nu,P)$ is the Russo--Van Roy information ratio.
\end{remark}

\section{Affine estimators and the open question}

The unrestricted optimizer $G^\nu$ can be computationally and conceptually unwieldy (it depends on the full posterior mean map $(A,Y)\mapsto \po$).
A natural restriction is to \emph{affine-in-$Y$} estimators.

\begin{definition}[Affine estimator class]
    Let $T_0,T_1:\cA\to\Rd$. Define
    \[
        E_T(A,Y)=T_1(A)\,Y+T_0(A).
    \]
    Write $\hat\ell = E_T(A,Y)$ and define the associated mirror update random variable
    \begin{equation}
        \label{eq:tildeA}
        \tilde A \;=\; \nabla F^\star\!\big(\nabla F(a)-\eta \hat\ell\big).
    \end{equation}
\end{definition}

\begin{lemma}[Value under an optimal affine estimator]
    \label{lem:affine-value}
    Fix $(a,\eta,P,\nu)$. Let $T^\nu$ be a minimizer of $\Lambda_{a,\eta}(P,E_T\Vert\nu)$ over affine estimators.
    With $\hat\ell=E_{T^\nu}(A,Y)$ and $\tilde A$ from \eqref{eq:tildeA}, one has
    \begin{equation}
        \label{eq:Lambda-affine}
        \Lambda_{a,\eta}(P,E_{T^\nu}\Vert\nu)
        \;=\;
        \frac{1}{\eta}\,\Delta(\nu,P)
        \;-\;
        \frac{1}{\eta^2}\,\E\!\big[D_F(\tilde A,a)\big].
    \end{equation}
    Moreover, the first-order conditions can be written as the conditional moment-matching relations
    \[
        \E[\tilde A \mid A]=\E[A_\star\mid A],\qquad
        \E[Y\,\tilde A \mid A]=\E[Y\,A_\star\mid A].
    \]
    (Since $A\sim P$ is independent of $(A_\star,\ell)\sim\nu$, $\E[A_\star\mid A]=\E[A_\star]=\pr$.)
\end{lemma}

The unrestricted and affine values are thus identical except for the substitution
\[
    \po=\E[A_\star\mid A,Y]
    \quad\leadsto\quad
    \tilde A=\nabla F^\star(\nabla F(a)-\eta E_T(A,Y)).
\]

A useful way to phrase the research target is via the (nonnegative) gap between the optimal affine value and the unrestricted value.

\begin{lemma}[Gap identity: ``affine is enough'' iff $\tilde A=\po$]
    \label{lem:gap}
    Under the affine first-order optimality condition, one can derive the decomposition
    \[
        \E[D_F(\po,a)] \;=\; \E[D_F(\po,\tilde A)] + \E[D_F(\tilde A,a)].
    \]
    Consequently,
    \[
        \Lambda_{a,\eta}(P,E_{T^\nu}\Vert\nu)-\Lambda_{a,\eta}(P,G^\nu\Vert\nu)
        \;=\;
        \frac{1}{\eta^2}\,\E[D_F(\po,\tilde A)]
        \;\ge\;0,
    \]
    with equality if and only if $\tilde A=\po$ almost surely (equivalently, $\E[D_F(\po,\tilde A)]=0$).
\end{lemma}

\begin{conjecture}[Affine estimators suffice for the information-ratio upper bound]
    In the minimax game defining $\Lambda^\star_\eta$ (or in the relevant saddle-point of Lemma~\ref{lem:minimax}),
    restricting the estimator class from all measurable $E:\cA\times\R\to\Rd$ to affine-in-$Y$ estimators does not increase the value.
    Equivalently, at the worst-case $\nu$ (and corresponding optimal $P$), there exists an affine estimator such that $\tilde A=\po$ a.s.,
    so the gap in Lemma~\ref{lem:gap} is zero and the information-ratio bound of Theorem~\ref{thm:info-ratio} is achieved within the affine class.
\end{conjecture}

\begin{remark}[What needs to be shown (minimal checklist)]
    A proof that ``affine is enough'' can be organized as one of the following equivalent targets:
    \begin{itemize}
        \item Show that at an optimal (or worst-case) $\nu$, the posterior mean map $(A,Y)\mapsto \po$ is affine in $Y$ for each fixed $A$.
        \item Show existence of affine $E_T$ such that the induced mirror update $\tilde A$ in \eqref{eq:tildeA} equals $\po$ a.s.
        \item Show the minimax value is unchanged after restricting estimators, i.e.
              \[
                  \sup_{\nu}\inf_{P}\inf_{E\ \mathrm{affine}}\Lambda_{a,\eta}(P,E\Vert\nu)
                  \;=\;
                  \sup_{\nu}\inf_{P}\inf_{E\ \mathrm{all}}\Lambda_{a,\eta}(P,E\Vert\nu),
              \]
              and then use Lemma~\ref{lem:gap} to conclude $\E[D_F(\po,\tilde A)]=0$ at the saddle point.
    \end{itemize}
\end{remark}

\end{document}
